\section{Problem Formulation}
\label{sec:problem}

\subsection{Fixed Fleet and Terminal Dispatch}
\label{sec:problem_corridor}

A bidirectional corridor has a physical fleet of $N_{\mathrm{fleet}}$ buses, capacity $C_{\mathrm{bus}}$, and a roster of $M_d$ trips in direction $d\in\{+,-\}$. Trips depart from one terminal and return to the other terminal's vehicle pool. The planner may revise uncommitted departure times $\tau_{d,m}^{\mathrm{sched}}$ and target headways $H_{d,m}$, but does not add or remove trips. The final planned departure in each direction is no earlier than the original final departure and the demand cutoff $T_{\mathrm{end}}$. We distinguish this planning constraint from completing the roster in the realized simulation.

A trip becomes eligible for dispatch when its revised scheduled time and minimum terminal separation permit it. It then waits if no vehicle is ready. The reported readiness delay is $D_{d,m}^{\mathrm{ready}}=\tau_{d,m}^{\mathrm{act}}-\tau_{d,m}^{\mathrm{elig}}$, where $\tau_{d,m}^{\mathrm{elig}}$ is the first eligible time. This is not delay from the original public timetable and excludes waiting before first eligibility. At intermediate stops a bus boards up to its remaining capacity. Its return time includes link travel, passenger-service dwell and discretionary holding:
\begin{equation}
  \tau^{\mathrm{ret}}_{v,m} = \tau^{\mathrm{act}}_{d,m} + T^{\mathrm{link}}_{v,m}
      + T^{\mathrm{dwell}}_{v,m} + \sum_s u_{v,m,s}.
  \label{eq:return_time}
\end{equation}
The simulator requires the vehicle to be available at the receiving terminal before reuse. The 20~s buffer used later in the ledger is a planning allowance, not a mandatory physical layover.

\subsection{Passenger and Regularity Metrics}
\label{sec:problem_metrics}

For the $|\mathcal{P}|$ passengers generated by the demand cutoff, restricted mean journey exposure $J_{\mathrm{pax}}$ is total waiting plus onboard person-time divided by $|\mathcal{P}|$. A passenger still waiting or riding at the evaluation cutoff contributes time accrued through that cutoff. We also report the fraction not served. The means do not condition on completing a journey.

For recorded nonterminal \emph{arrival} headways $h$, we report raw $\mathrm{CV}^2_{\mathrm{raw}}=\mathrm{Var}(h)/\mathbb{E}[h]^2$. Let $\mathcal{E}$ contain the observed headways and their active target headways $H$. Plan-normalized headway error and the excess-waiting-time (EWT) proxy are:
\begin{equation}
\begin{aligned}
  \mathrm{CV}^2_{\mathrm{plan}} &=
     \frac{1}{|\mathcal{E}|}\sum_{(h,H)\in\mathcal E}\left(\frac{h-H}{H}\right)^2,\\
  \mathrm{EWT} &= \frac{1}{|\mathcal{S}|}\sum_{(s,d)\in\mathcal{S}}
  \left[\frac{\sum_j h_{s,d,j}^2}{2\sum_j h_{s,d,j}} -
        \frac{\sum_j H_{s,d,j}^2}{2\sum_j H_{s,d,j}}\right],
\end{aligned}
\label{eq:plan_cv2_ewt}
\end{equation}
where $\mathcal{S}$ is the set of stop-direction groups with valid headways. EWT is relative to each controller's own plan, not a shared public timetable or a direct measure of actual passenger waits. The raw CV$^2$ does not remove intentional headway changes. Neither metric substitutes for the direct passenger outcome.

\subsection{A Plan-Consistent Two-Gap Holding Target}

\begin{proposition}[Plan-Consistent Non-Stationary Bilateral Holding Target]
\label{prop:nonstationary_target}
Suppose vehicle $v$ serving trip $(d,m)$ with planned headway $H_{\mathrm{own}} > 0$ completes boarding at stop $s$ with causal forward departure headway $g^{\mathrm{fwd}} \ge 0$ and same-time causal follower departure gap $g^{\mathrm{bwd}} \ge 0$, where the trailing vehicle serves trip $(d,m+1)$ with planned headway $H_{\mathrm{fol}} > 0$. Applying holding $u \in [0,\bar u]$ with $\bar u>0$ shifts the local gaps to $(g^{\mathrm{fwd}} + u,\, g^{\mathrm{bwd}} - u)$. Minimizing the local plan-normalized regularity loss
\begin{equation}
  \mathcal{L}_{\mathrm{reg}}(u) = \frac{1}{2}\left[\left(\frac{g^{\mathrm{fwd}} + u - H_{\mathrm{own}}}{H_{\mathrm{own}}}\right)^2 + \left(\frac{g^{\mathrm{bwd}} - u - H_{\mathrm{fol}}}{H_{\mathrm{fol}}}\right)^2\right]
  \label{eq:local_reg_loss}
\end{equation}
over $u \in [0, \bar{u}]$ yields the unique closed-form target
\begin{equation}
  u^\star = \mathrm{clip}\!\left(
    \frac{(g^{\mathrm{bwd}} - H_{\mathrm{fol}}) H_{\mathrm{fol}}^{-2} - (g^{\mathrm{fwd}} - H_{\mathrm{own}}) H_{\mathrm{own}}^{-2}}{H_{\mathrm{own}}^{-2} + H_{\mathrm{fol}}^{-2}},
    \, 0,\, \bar{u}
  \right).
  \label{eq:u_star_nonstationary}
\end{equation}
In particular, when $H_{\mathrm{own}} = H_{\mathrm{fol}}$, Eq.~\eqref{eq:u_star_nonstationary} reduces to the classical bilateral rule $u_{\mathrm{stat}}^\star = \mathrm{clip}(\frac{1}{2}(g^{\mathrm{bwd}} - g^{\mathrm{fwd}}), 0, \bar{u})$. However, when vehicles follow a non-stationary schedule perfectly ($g^{\mathrm{fwd}} = H_{\mathrm{own}}$, $g^{\mathrm{bwd}} = H_{\mathrm{fol}}$) with $H_{\mathrm{fol}} > H_{\mathrm{own}}$ (peak-to-off-peak transition), Eq.~\eqref{eq:u_star_nonstationary} yields $u^\star = 0$, whereas the stationary rule injects spurious holding $u_{\mathrm{stat}}^\star = \frac{1}{2}(H_{\mathrm{fol}} - H_{\mathrm{own}}) > 0$.
\end{proposition}

\begin{proof}
The objective $\mathcal{L}_{\mathrm{reg}}(u)$ in Eq.~\eqref{eq:local_reg_loss} is strictly convex in $u$ with second derivative $\mathcal{L}_{\mathrm{reg}}''(u) = H_{\mathrm{own}}^{-2} + H_{\mathrm{fol}}^{-2} > 0$. Setting the first derivative to zero gives
\begin{align*}
  &\frac{g^{\mathrm{fwd}} + u - H_{\mathrm{own}}}{H_{\mathrm{own}}^2} - \frac{g^{\mathrm{bwd}} - u - H_{\mathrm{fol}}}{H_{\mathrm{fol}}^2} = 0 \\
  &\iff u\bigl(H_{\mathrm{own}}^{-2} + H_{\mathrm{fol}}^{-2}\bigr) = (g^{\mathrm{bwd}} - H_{\mathrm{fol}}) H_{\mathrm{fol}}^{-2} - (g^{\mathrm{fwd}} - H_{\mathrm{own}}) H_{\mathrm{own}}^{-2}.
\end{align*}
Projecting the unconstrained minimizer onto the compact interval $[0, \bar{u}]$ yields Eq.~\eqref{eq:u_star_nonstationary}. Substituting $H_{\mathrm{own}} = H_{\mathrm{fol}} = H$ cancels $H^{-2}$ and cancels $\pm H$ in the numerator, yielding $\frac{1}{2}(g^{\mathrm{bwd}} - g^{\mathrm{fwd}})$. Substituting $g^{\mathrm{fwd}} = H_{\mathrm{own}}$ and $g^{\mathrm{bwd}} = H_{\mathrm{fol}}$ makes the numerator identically zero, so $u^\star = 0$, whereas $\frac{1}{2}(g^{\mathrm{bwd}} - g^{\mathrm{fwd}}) = \frac{1}{2}(H_{\mathrm{fol}} - H_{\mathrm{own}}) > 0$ whenever $H_{\mathrm{fol}} > H_{\mathrm{own}}$.
\end{proof}

\subsection{Local Waiting Savings and the Limits of Turnaround Pricing}
\label{sec:problem_osuna}

Under a constant arrival rate $\Lambda$ (passengers per second), no capacity limit and rigid shifts on a downstream two-gap segment, the expected waiting saved by holding $u$ is
\begin{equation}
  \Delta W(u)=\Lambda u(g^{\mathrm{bwd}}-g^{\mathrm{fwd}}-u).
  \label{eq:osuna_turnaround_delta}
\end{equation}
This follows by subtracting the squared-gap waiting cost after holding from that before holding \citep{osuna1972control}. The local onboard cost is $L^{\mathrm{on}}u$. When targets differ, the controller instead uses a plan-aware \emph{proxy} for the first term:
\begin{equation}
  G(u)=\widehat{\Lambda}^{\mathrm{ds}}\frac{H_{\mathrm{own}}+H_{\mathrm{fol}}}{2}u
  \left[\frac{g^{\mathrm{bwd}}}{H_{\mathrm{fol}}}-\frac{g^{\mathrm{fwd}}}{H_{\mathrm{own}}}
  -\frac{u}{2}\!\left(\frac{1}{H_{\mathrm{own}}}+\frac{1}{H_{\mathrm{fol}}}\right)\right].
  \label{eq:normalized_hold_proxy}
\end{equation}
It reduces to Eq.~\eqref{eq:osuna_turnaround_delta} for equal headway targets with $\Lambda=\widehat{\Lambda}^{\mathrm{ds}}$. With unequal targets, capacity constraints or state-dependent dwell, it is not the causal change in system-wide waiting.

Turnaround effects depend on future departures. For illustration, if the same vehicle is required for the next dispatch, a hold beyond its actual slack can delay that dispatch. Whether the resulting net waiting cost is large also depends on the next gap. With constant arrivals, initially equal gaps $H$, and a single dispatch delayed by $\delta<H$ while its neighbors stay fixed, gaps become $H+\delta$ and $H-\delta$, and the added waiting across both gaps is $\Lambda\delta^2$. Other ready vehicles may absorb the delay entirely. We therefore use predicted slack to limit holds, and measure passenger and vehicle consequences in simulation rather than assuming a lower bound on spillover.
